\textbf{2.\ (Moving the knight).}

We number the board by coordinates
\[
a1=(1,1),\qquad h8=(8,8).
\]
The knight starts at
\[
h8=(8,8).
\]

A legal move is one of
\[
(x,y)\mapsto (x+1,y-2),
\]
\[
(x,y)\mapsto (x-1,y-2),
\]
\[
(x,y)\mapsto (x-2,y+1),
\]
\[
(x,y)\mapsto (x-2,y-1)
\]
provided the new square is still on the board.
winning if the player to move from it can move to a losing position.

\subsection*{Losing positions on the \(8\times 8\) board}

In coordinates, these are the \(2\times2\) blocks
\[
\{1,2\}\times\{1,2\},
\]
\[
\{1,2\}\times\{5,6\},
\]
\[
\{5,6\}\times\{1,2\},
\]
\[
\{5,6\}\times\{5,6\},
\]
together with the special square
\[
(8,8).
\]

We can see that losing positions satisfy
\[
x,y\equiv 1 \text{ or } 2 \pmod 4
\]

On the \(8\times 8\) board, the losing positions can be displayed as follows:

\[
\begin{array}{c|cccccccc}
& 1 & 2 & 3 & 4 & 5 & 6 & 7 & 8\\
\hline
8 & . & . & . & . & . & . & . & L\\
7 & . & . & . & . & . & . & . & .\\
6 & L & L & . & . & L & L & . & .\\
5 & L & L & . & . & L & L & . & .\\
4 & . & . & . & . & . & . & . & .\\
3 & . & . & . & . & . & . & . & .\\
2 & L & L & . & . & L & L & . & .\\
1 & L & L & . & . & L & L & . & .
\end{array}
\]

\subsection*{Why these positions are losing}

Let
\[
R=\{(x,y): x,y\equiv 1 \text{ or }2 \pmod 4\}.
\]

These positions form repeated \(2\times2\) blocks. A knight move changes one
coordinate by \(1\) and the other coordinate by \(2\). Therefore a move from a
position in \(R\) always leaves \(R\). So no position in \(R\) has a move to
another position in \(R\).

Conversely, by backward induction, every legal move out of one of these blocks
leads to a position from which the next player can move back into a smaller
\(R\)-position. Hence the positions in \(R\) are losing.

Now check the starting square:
\[
h8=(8,8)
\]
Its only legal moves are
\[
(8,8)\mapsto (7,6)
\]
and
\[
(8,8)\mapsto (6,7)
\]

Both of these positions are winning, because both can move to
\[
(5,5)=e5,
\]
which is losing:
\[
(7,6)\mapsto (5,5),
\]
\[
(6,7)\mapsto (5,5)
\]

Therefore every legal move from \(h8\) goes to a winning position. Hence
\[
h8
\]
is losing.

Thus the first player loses and the second player wins.

% \subsection*{General \(n\times n\) board}

% For the same game on an \(n\times n\) board, the losing block pattern repeats
% modulo \(4\): positions with
% \[
% x,y\equiv 1 \text{ or }2 \pmod 4
% \]
% are losing.

% Starting from \((n,n)\), if \(n\equiv1\) or \(2\pmod4\), the initial square is
% already losing. If \(n\equiv3\pmod4\), the first player moves
% \[
% (n,n)\mapsto(n-1,n-2),
% \]
% which has residue \((2,1)\) and is losing. If \(n\equiv0\pmod4\), the only
% first moves go to \((n-1,n-2)\) or \((n-2,n-1)\), and each can be answered by
% moving to \((n-3,n-3)\equiv(1,1)\pmod4\), a losing square.

% Thus, on an \(n\times n\) board starting from the top-right corner,

% \[
% \boxed{\text{the first player wins if and only if } n\equiv 3\pmod4.}
% \]

% For the usual chessboard,
% \[
% 8\equiv0\pmod4,
% \]
% so the first player loses and the second player wins.

\section*{Moving the knight: general solution}

Consider an \(n\times n\) board with coordinates
\[
(1,1), (1,2), \ldots, (n,n),
\]
where the knight starts at the top-right corner
\[
(n,n).
\]

The legal moves are the same as before whenever the new square is still on the board.


\subsection*{Losing blocks}

The basic losing positions are the \(2\times2\) blocks satisfying
\[
x,y\equiv 1 \text{ or } 2 \pmod 4.
\]


A move from such a square always leaves this set. Moreover, after any legal move from such a square, the next player can move back to a smaller square of the same type. Therefore these positions are losing positions.

\subsection*{Starting from the corner}

Now consider the starting square
\[
(n,n).
\]

There are four cases.

\subsubsection*{Case 1: \(n\equiv1\pmod4\)}

Then
\[
(n,n)\equiv(1,1)\pmod4,
\]
so \((n,n)\) lies in a losing block. Hence the first player loses.

\subsubsection*{Case 2: \(n\equiv2\pmod4\)}

Then
\[
(n,n)\equiv(2,2)\pmod4,
\]
so \((n,n)\) also lies in a losing block. Hence the first player loses.

\subsubsection*{Case 3: \(n\equiv3\pmod4\)}

Then the first player can move
\[
(n,n)\mapsto(n-1,n-2)
\]
Modulo \(4\), this is
\[
(3,3)\mapsto(2,1)
\]
Since \((2,1)\) is a losing-type residue, the first player can move to a losing position. Hence the first player wins.

\subsubsection*{Case 4: \(n\equiv0\pmod4\)}

The square \((n,n)\) is not itself in the losing block pattern. However, because it is the top-right corner, it has only two legal moves:
\[
(n,n)\mapsto(n-1,n-2),
\]
\[
(n,n)\mapsto(n-2,n-1)
\]

Both resulting positions are winning, because from either one the next player can move to
\[
(n-3,n-3).
\]

Indeed,
\[
(n-1,n-2)\mapsto(n-3,n-3),
\]
and
\[
(n-2,n-1)\mapsto(n-3,n-3)
\]

Since
\[
n\equiv0\pmod4,
\]
we have
\[
n-3\equiv1\pmod4
\]
Therefore
\[
(n-3,n-3)\equiv(1,1)\pmod4,
\]
which is a losing block position.

Thus every legal move from \((n,n)\) gives the opponent a move to a losing position. Hence \((n,n)\) is losing.

\subsection*{Conclusion}

Therefore, on an \(n\times n\) board starting from the top-right corner,
\[
\text{the first player wins if and only if } n\equiv3\pmod4
\]

For the usual chessboard,
\[
8\equiv0\pmod4,
\]
so the first player loses and the second player wins.
